
\kafkasection{A spherical fireball with mass loading $M$ expands into a constant density medium with one hydrogen atom per $cm^{-3}$. At which radius has it swept up a mass $M$? How does this radius change if the fireball is instead collimated into two cones with opening angle $\theta_j$?}
Since mass is just density times a volume, we will use the volume of a sphere. If the fireball has some initial radius $r_0$ the area needed to sweep out a mass $M$ is then:
\begin{equation}
    M = \rho(V_{outer} - V_{r_0}) = \rho\frac{4\pi}{3}\ab(r^3 - r_0^3)
\end{equation}
Of course, this just goes to the volume of the outer sphere if we have $r_0 = 0$. Rearranging for the radius is then:
\begin{equation}
    r = \ab(\frac{3M}{4\pi\rho})^{1/3}
\end{equation}
For when we extend out into 2 cones. Simply replace the the volume of the sphere with that of 2 cones.
\begin{equation}
    M = 2\rho\ab(\frac{\pi\tan^2\theta_j r^3}{3})
\end{equation}
\begin{equation}
    r = \ab(\frac{3M}{2\pi\rho\tan^2\theta_j})^{1/3}
\end{equation}
A similar formula where the demoninator varies as $\theta_j$. 
